\begin{table}[H]
\centering
\caption{Standardized Effect Sizes for Main Outcomes}
\label{tab:sde}
\begin{threeparttable}
\begin{adjustbox}{max width=\textwidth}
\begin{tabular}{llccccc}
\toprule
Outcome & Specification & $\hat{\beta}$ & SD($Y$) & SDE & SE(SDE) & Classification \\
\midrule
\multicolumn{7}{l}{\textit{Panel A: Pooled}} \\
Monthly hours & Baseline DiD & -0.212 & 15.48 & -0.0137 & 0.0689 & Small negative \\
Monthly days & Baseline DiD & -0.048 & 1.19 & -0.0406 & 0.0677 & Small negative \\
\midrule
\multicolumn{7}{l}{\textit{Panel B: Heterogeneous (by sex)}} \\
Monthly hours & Male workers & -0.181 & 12.83 & -0.0141 & 0.0724 & Small negative \\
Monthly hours & Female workers & 1.035 & 16.48 & 0.0628 & 0.0487 & Moderate positive \\
\bottomrule
\end{tabular}
\end{adjustbox}
\begin{tablenotes}[flushleft]
\scriptsize\sloppy
\item \textit{Notes:} \textbf{Country:} Japan. \textbf{Research question:} Does the application of statutory overtime caps to previously exempted industries reduce worker hours and days worked? \textbf{Policy mechanism:} The 2018 Work Style Reform Act imposed binding monthly and annual overtime caps on most Japanese industries from April 2019, but granted five-year exemptions to transport, construction, and healthcare due to their ``distinctive'' labor demands. When exemptions expired in April 2024, these industries faced the same caps as all other sectors, creating a compliance cliff where firms had to immediately limit previously unconstrained overtime scheduling. \textbf{Outcome definition:} Average monthly hours worked per employed person, measured by the Labour Force Survey (household-based) at the industry-month level. \textbf{Treatment:} Binary---industry exempt from overtime caps until April 2024 vs.~already subject to caps since April 2019. \textbf{Data:} Labour Force Survey via e-Stat API, April 2017--January 2026, industry-month level, 19 major industries (excluding aggregates). \textbf{Method:} Two-way fixed effects DiD with industry and year-month fixed effects, standard errors clustered at the industry level, supplemented by randomization inference (2,000 permutations). \textbf{Sample:} All employed persons across 19 major industry divisions; 3 exempt industries (Construction, Transport/Postal, Medical/Welfare) and 16 non-exempt industries. SDE $= \hat{\beta} / \text{SD}(Y)$ where SD($Y$) is the pre-treatment standard deviation. Classification refers to magnitude, not statistical significance: Large ($|$SDE$| > 0.15$), Moderate ($0.05$--$0.15$), Small ($0.005$--$0.05$), Null ($< 0.005$).
\end{tablenotes}
\end{threeparttable}
\end{table}
